\documentclass[UTF8]{ctexart}
\usepackage[a4paper,margin=2.4cm]{geometry}
\usepackage{amsmath,amssymb,mathtools}
\usepackage{graphicx}
\usepackage{booktabs}
\usepackage{hyperref}

\title{第一题：LIF 模型的更新过程推导与数值模拟}
\author{}
\date{\today}

\begin{document}
\emergencystretch=2em
\maketitle

\section{模型与记号}

题目给出的 LIF 模型为
\begin{align}
  \frac{dV}{dt} &= -g_L V + g_{\mathrm{syn}} I_{\mathrm{syn}}, \label{eq:lif_v}\\
  \tau \frac{d I_{\mathrm{syn}}}{dt}
  &= - I_{\mathrm{syn}}+\mu+\sigma \frac{dW}{dt}, \label{eq:lif_i}
\end{align}
其中 $W$ 为标准 Wiener 过程。膜电位达到阈值 $V_{\mathrm{th}}$ 时产生一次动作电位，
随后经历绝对不应期 $\tau_{\mathrm{ref}}$。题目未给出重置电位，本文取
$V_{\mathrm{r}}=0$。连续两次发放之间的间隔记为
\[
T=\tau_{\mathrm{ref}}+\tau_h,
\]
其中 $\tau_h$ 是从重置电位 $V_{\mathrm r}$ 首次到达阈值 $V_{\mathrm{th}}$ 的首达时间。
若每次发放后系统按同一规则重置，则 $T_1,T_2,\ldots$ 可看作独立同分布的更新间隔，
从而放电计数过程
\[
N(t)=\max\{n:S_n\le t\},\qquad S_n=T_1+\cdots+T_n
\]
构成普通更新过程。

\section{第（1）问：ISI 矩与更新函数的推导}

\subsection{将突触电流近似为有效扩散输入}

由式 \eqref{eq:lif_i} 可写成
\[
dI_{\mathrm{syn}}
=\frac{\mu-I_{\mathrm{syn}}}{\tau}\,dt+\frac{\sigma}{\tau}\,dW_t.
\]
这是均值为 $\mu$ 的 Ornstein--Uhlenbeck 过程，其稳态协方差为
\[
\operatorname{Cov}\bigl(I_{\mathrm{syn}}(t),I_{\mathrm{syn}}(t+s)\bigr)
=\frac{\sigma^2}{2\tau}e^{-|s|/\tau}.
\]
当突触相关时间 $\tau$ 相对放电时间尺度较小时，有
\[
\operatorname{Var}\left(\int_0^t (I_{\mathrm{syn}}(s)-\mu)\,ds\right)
\approx t\cdot 2\int_0^\infty \frac{\sigma^2}{2\tau}e^{-s/\tau}\,ds
=\sigma^2 t.
\]
因此可用有效一维扩散近似
\begin{equation}
  dV_t=(-g_LV_t+A)\,dt+B\,dB_t,\qquad
  A=g_{\mathrm{syn}}\mu,\quad B=g_{\mathrm{syn}}\sigma,
  \label{eq:effective_diffusion}
\end{equation}
来描述膜电位的首达时间。这里 $B_t$ 是另一个标准 Brownian motion。

\subsection{首达时间一阶矩和二阶矩}

对扩散过程 \eqref{eq:effective_diffusion}，其生成元为
\[
\mathcal L f(x)=(-g_Lx+A)f'(x)+\frac{B^2}{2}f''(x).
\]
设
\[
u_n(x)=\mathbb E_x[\tau_h^n],\qquad u_0(x)=1.
\]
由 Dynkin 公式或后向 Kolmogorov 方程，首达时间矩满足递推边值问题
\begin{equation}
\mathcal L u_n(x)=-n u_{n-1}(x),\qquad x<V_{\mathrm{th}},
\label{eq:moment_ode}
\end{equation}
并具有边界条件
\[
u_n(V_{\mathrm{th}})=0,\qquad u_n(x)\text{ 在 }x\to-\infty\text{ 时有界}.
\]
具体地，一阶矩满足
\begin{equation}
(-g_Lx+A)u_1'(x)+\frac{B^2}{2}u_1''(x)=-1,
\qquad u_1(V_{\mathrm{th}})=0,
\label{eq:first_moment_ode}
\end{equation}
二阶矩满足
\begin{equation}
(-g_Lx+A)u_2'(x)+\frac{B^2}{2}u_2''(x)=-2u_1(x),
\qquad u_2(V_{\mathrm{th}})=0.
\label{eq:second_moment_ode}
\end{equation}
令
\[
\Phi(x)=\int^x \frac{2(-g_Ly+A)}{B^2}\,dy
=\frac{2Ax-g_Lx^2}{B^2}.
\]
将 \eqref{eq:moment_ode} 视为关于 $u_n'(x)$ 的一阶线性方程，可得递推积分表达式
\begin{equation}
u_n(x)=
n\int_x^{V_{\mathrm{th}}}
e^{-\Phi(y)}
\int_{-\infty}^{y}\frac{2}{B^2}e^{\Phi(z)}u_{n-1}(z)\,dz\,dy .
\label{eq:moment_integral}
\end{equation}
于是 ISI 的均值和方差为
\begin{align}
\mathbb E[T]
&=\tau_{\mathrm{ref}}+u_1(V_{\mathrm r}), \label{eq:isi_mean}\\
\operatorname{Var}(T)
&=u_2(V_{\mathrm r})-u_1(V_{\mathrm r})^2. \label{eq:isi_var}
\end{align}
注意不应期是确定常数，因此只改变均值，不改变方差。

\subsection{更新函数的期望表达式}

记 $F(t)=\mathbb P(T\le t)$，$F^{*n}$ 为 $F$ 的 $n$ 重卷积。由于
\[
N(t)\ge n \iff S_n=T_1+\cdots+T_n\le t,
\]
有
\[
\mathbb P(N(t)\ge n)=F^{*n}(t).
\]
利用非负整数随机变量恒等式
\[
N(t)=\sum_{n=1}^{\infty}\mathbf 1_{\{N(t)\ge n\}},
\]
两边取期望得到更新函数
\begin{equation}
M(t)=\mathbb E[N(t)]
=\sum_{n=1}^{\infty}F^{*n}(t).
\label{eq:renewal_mean_exact}
\end{equation}
若 $\widehat F(s)$ 表示 $F$ 的 Laplace--Stieltjes 变换，则
\[
\mathcal L\{F^{*n}\}(s)=\frac{\widehat F(s)^n}{s},
\]
因此
\begin{equation}
\widehat M(s)
=\frac{1}{s}\sum_{n=1}^{\infty}\widehat F(s)^n
=\frac{\widehat F(s)}{s(1-\widehat F(s))}.
\label{eq:renewal_laplace}
\end{equation}

\subsection{更新计数方差表达式}

同样利用整数变量恒等式
\[
N(t)^2=\sum_{n=1}^{N(t)}(2n-1)
=\sum_{n=1}^{\infty}(2n-1)\mathbf 1_{\{N(t)\ge n\}},
\]
取期望得
\begin{equation}
\mathbb E[N(t)^2]
=\sum_{n=1}^{\infty}(2n-1)F^{*n}(t).
\label{eq:renewal_second_exact}
\end{equation}
从而
\begin{equation}
\operatorname{Var}(N(t))
=\sum_{n=1}^{\infty}(2n-1)F^{*n}(t)-M(t)^2.
\label{eq:renewal_var_exact}
\end{equation}
其二阶矩的 Laplace 形式为
\begin{equation}
\mathcal L\{\mathbb E[N(t)^2]\}(s)
=\frac{1}{s}\sum_{n=1}^{\infty}(2n-1)\widehat F(s)^n
=\frac{\widehat F(s)(1+\widehat F(s))}
{s(1-\widehat F(s))^2}.
\label{eq:renewal_second_laplace}
\end{equation}
若只关心长时间近似，并令
\[
m_T=\mathbb E[T],\qquad v_T=\operatorname{Var}(T),
\]
则普通更新定理和更新中心极限定理给出
\begin{equation}
\mathbb E[N(t)]\sim \frac{t}{m_T},
\qquad
\operatorname{Var}(N(t))\sim \frac{v_T}{m_T^3}\,t,
\qquad t\to\infty.
\label{eq:renewal_asymptotic}
\end{equation}
这说明平均放电率为 $1/m_T$，而计数方差的增长率由 ISI 方差 $v_T$ 决定。

\section{第（2）问：给定目标矩时的参数构造}

给定目标
\[
\mathbb E[T]\approx m,\qquad \operatorname{Var}(T)\approx s^2,
\]
由 \eqref{eq:isi_mean}--\eqref{eq:isi_var} 可知，只需要调节有效漂移 $A$ 和有效扩散强度
$B$。因为
\[
A=g_{\mathrm{syn}}\mu,\qquad B=g_{\mathrm{syn}}\sigma,
\]
三元组 $(g_{\mathrm{syn}},\mu,\sigma)$ 并不唯一。本文取 $g_{\mathrm{syn}}=1$，
然后令
\[
\mu=A,\qquad \sigma=B.
\]
若希望固定其他正值 $g_{\mathrm{syn}}=c$，则只需取
\[
\mu=\frac{A}{c},\qquad \sigma=\frac{B}{c}.
\]

具体构造步骤如下。

\begin{enumerate}
\item 先扣除不应期，令目标首达时间均值为
\[
\bar m=m-\tau_{\mathrm{ref}}.
\]
必须有 $\bar m>0$。
\item 对给定 $B$，用有限差分求解 \eqref{eq:first_moment_ode}--\eqref{eq:second_moment_ode}。
边界上取 $u_n(V_{\mathrm{th}})=0$，左端取足够远的反射边界近似自然边界。
\item 对给定 $B$，用二分法寻找 $A$，使
\[
u_1(V_{\mathrm r};A,B)=\bar m.
\]
\item 再对 $B$ 做二分法，使
\[
u_2(V_{\mathrm r};A,B)-u_1(V_{\mathrm r};A,B)^2=s^2.
\]
\item 得到 $A,B$ 后构造
\[
(g_{\mathrm{syn}},\mu,\sigma)=\left(1,A,B\right).
\]
\end{enumerate}

这里还需要说明一个数值边界处理。理论推导中，后向方程的左边界条件写为
$x\to-\infty$ 时 $u_n(x)$ 有界，这对应扩散过程在远离阈值的左侧没有吸收边界。
但在 Python 中不能在无穷区间上直接求解有限差分方程，因此代码把计算区间截断为
\[
[x_{\min},V_{\mathrm{th}}],
\qquad
x_{\min}
=\min\left\{
V_{\mathrm r}-7\frac{B}{\sqrt{2g_L}}-1,\,-2
\right\}.
\]
其中 $B/\sqrt{2g_L}$ 是有效 OU 型膜电位扩散的稳态标准差尺度。这样选取
$x_{\min}$ 可以保证左边界远低于重置电位和主要概率质量区域。数值上在
$x_{\min}$ 处使用反射边界
\[
u_n(x_0)-u_n(x_1)=0,
\]
也就是近似令 $u_n'(x_{\min})=0$。这不是改变理论模型，而是用一个足够远的有限边界
近似理论中的自然边界条件；当 $x_{\min}$ 继续向左移动时，该截断误差会迅速减小。

数值搜索的初值可由无噪声确定性 LIF 给出。若 $B=0$，膜电位轨道为
\[
V(t)=\frac{A}{g_L}+\left(V_{\mathrm r}-\frac{A}{g_L}\right)e^{-g_Lt}.
\]
令 $V(\bar m)=V_{\mathrm{th}}$，得到
\[
A_0=
\frac{g_L\left(V_{\mathrm{th}}-V_{\mathrm r}e^{-g_L\bar m}\right)}
{1-e^{-g_L\bar m}}.
\]
这给出了漂移大小的合理初始范围；随后用后向方程精确匹配目标矩。

\section{Python 数值模拟}

程序文件为 \texttt{lif\_renewal\_sim.py}。模拟使用原始二维模型
\begin{align*}
dI_{\mathrm{syn}}
&=\frac{\mu-I_{\mathrm{syn}}}{\tau}\,dt+\frac{\sigma}{\tau}\,dW_t,\\
dV_t&=(-g_LV_t+g_{\mathrm{syn}}I_{\mathrm{syn}})\,dt,
\end{align*}
而不是只模拟一维扩散近似。每个 ISI 从 $V_{\mathrm r}=0$ 开始，$I_{\mathrm{syn}}(0)$
从 OU 稳态分布
\[
I_{\mathrm{syn}}(0)\sim \mathcal N\left(\mu,\frac{\sigma^2}{2\tau}\right)
\]
抽样，用 Euler--Maruyama 格式推进；到达阈值后用线性插值估计穿越时间，并加上
$\tau_{\mathrm{ref}}$ 得到完整 ISI。

本次默认目标和常数为
\[
m=1.0,\qquad s^2=0.08,\qquad
g_L=0.2,\quad \tau=0.02,\quad
V_{\mathrm r}=0,\quad V_{\mathrm{th}}=1,\quad \tau_{\mathrm{ref}}=0.05.
\]
运行命令为
\[
\texttt{python lif\_renewal\_sim.py}.
\]

\section{代码结果与解释}

程序构造出的参数为
\begin{center}
\begin{tabular}{ll}
\toprule
参数 & 数值 \\
\midrule
$g_{\mathrm{syn}}$ & $1.000000$ \\
$\mu$ & $1.146884$ \\
$\sigma$ & $0.306647$ \\
$A=g_{\mathrm{syn}}\mu$ & $1.146884$ \\
$B=g_{\mathrm{syn}}\sigma$ & $0.306647$ \\
\bottomrule
\end{tabular}
\end{center}

理论构造与原始模型仿真对比如下：
\begin{center}
\begin{tabular}{lcc}
\toprule
来源 & $\mathbb E[T]$ & $\operatorname{Var}(T)$ \\
\midrule
目标值 & $1.000000$ & $0.080000$ \\
扩散近似后向方程 & $1.000000$ & $0.080000$ \\
原始二维 SDE 仿真 & $1.032638$ & $0.079216$ \\
\bottomrule
\end{tabular}
\end{center}
仿真使用 $8000$ 个独立 ISI，时间步长为 $5\times 10^{-4}$。相对误差为
\[
\frac{\widehat{\mathbb E[T]}-m}{m}=3.26\%,\qquad
\frac{\widehat{\operatorname{Var}(T)}-s^2}{s^2}=-0.98\%.
\]
均值有小幅偏大，主要来自三类误差。第一，参数构造基于短相关时间极限下的一维有效扩散近似，
而验证时使用的是带有限相关时间 $\tau=0.02$ 的原始 OU 突触电流。更具体地，
\[
\operatorname{Var}\left(\int_0^t(I_{\mathrm{syn}}(s)-\mu)\,ds\right)
=\sigma^2\left[t-\tau\left(1-e^{-t/\tau}\right)\right],
\]
只有在 $\tau\to0$ 或 $t\gg \tau$ 的极限下才收敛到有效扩散模型使用的
$\sigma^2t$。因此随着 $\tau\to0$，原始二维模型会更接近一维扩散近似，均值误差也应进一步收敛。
第二，阈值穿越由 Euler--Maruyama 离散步长 $dt=5\times10^{-4}$ 和线性插值近似，
仍会留下小的离散化误差。第三，Monte Carlo 样本量有限；本次 $8000$ 个 ISI 下，
均值标准误约为 $\sqrt{0.079216/8000}\approx0.0031$，明显小于 $3.26\%$ 的系统偏差，
说明有限相关时间近似是主要来源，而方差的 $-0.98\%$ 偏差已在采样波动量级内。
总体上，构造参数能较好地达到给定 ISI 方差，并使均值保持在目标附近。

\begin{figure}[htbp]
\centering
\includegraphics[width=0.82\linewidth]{results/isi_histogram.png}
\caption{模拟得到的 ISI 分布直方图}
\label{fig:isi_hist}
\end{figure}

图 \ref{fig:isi_hist} 给出了 $8000$ 个 ISI 样本的直方图。红色竖线是目标均值
$m=1.0$，绿色竖线是原始二维 SDE 仿真得到的样本均值 $1.032638$。两条竖线接近，
说明参数构造达到了题目要求的“等于或者近似”目标。直方图略向右偏，这是首达时间分布的典型形态：
大多数轨道在平均时间附近穿越阈值，但少数受到负向噪声影响的轨道会显著延迟。

\begin{figure}[htbp]
\centering
\includegraphics[width=0.86\linewidth]{results/renewal_counts.png}
\caption{更新计数 $N(t)$ 的经验均值、方差与渐近公式对比}
\label{fig:renewal_counts}
\end{figure}

图 \ref{fig:renewal_counts} 用模拟得到的 ISI 样本重抽样生成多个更新过程，计算不同时间 $t$
下 $N(t)$ 的经验均值与经验方差。红色实线为经验 $\mathbb E[N(t)]$，红色虚线为
$t/\widehat m_T$；蓝色实线为经验 $\operatorname{Var}(N(t))$，蓝色虚线为
$\widehat v_Tt/\widehat m_T^3$。可以看到，随着时间增大，经验均值和方差都沿着更新理论给出的
渐近直线增长，验证了第（1）问中的更新函数推导。

\section{结论}

第（1）问中，ISI 可表示为确定不应期加扩散过程首达时间，其矩由后向 Kolmogorov 方程
\eqref{eq:moment_ode} 递推给出；更新计数的均值和方差则由 $F^{*n}$ 的卷积级数给出，
长时间下满足 $\mathbb E[N(t)]\sim t/m_T$ 与
$\operatorname{Var}(N(t))\sim v_Tt/m_T^3$。第（2）问中，利用
$A=g_{\mathrm{syn}}\mu$、$B=g_{\mathrm{syn}}\sigma$ 的非唯一性，可先求有效参数
$A,B$，再构造 $(g_{\mathrm{syn}},\mu,\sigma)$。数值模拟表明，在默认目标
$m=1.0,s^2=0.08$ 下，构造参数使原始模型的 ISI 均值和方差与目标值较为接近。

\end{document}
